Las herramientas encargadas de llevar a cabo la toma de métricas de consumo de las cargas de codificación desplegadas sobre contenedores de Docker tendrían que estar enfocadas a entornos de altas prestaciones compuestos por múltiples nodos. Además, debería ser capaz de tomar métricas de consumo de recursos hardware implicados, como la GPU, y poder diferenciar qué porcentaje de recursos ha sido utilizado por el contenedor.

Teniendo en cuenta estas necesidades, se ha realizado una búsqueda entre diversos repositorios, fuentes académicas y papers, como PAPER-KEPLER, y se ha llevado a cabo una selección de las herramientas que más se mencionan, tanto de forma directa como en la bibliografía. Estas dos herramientas han sido: Kepler \cite{kepler} y Scaphandre \cite{scaphandre}.

Además, también se han proporcionado dos herramientas por parte de Juan Gutiérrez Aguado, uno de los directores del proyecto Eco-Stream: EnergyMeter, empleada como referencia de consumo energetico del sistema basado en las metricas del sistema obtenidas por RAPL; Y Monitor, una herramienta desarrollada por Juan Gutiérrez Aguado empleada como alternativa a las herramientas anteriores, la cual permite medir el consumo de una carga de trabajo desplegada sobre uno o varios contendores de forma simultanea y muy precisa.\

%TABLA DE HERRAMIENTAS CAPACIDAD - NO SE SI FUNCIONA
\begin{table}[]
     \begin{tabular}{|c|c|c|c|c|}
     \hline
     \rowcolor[HTML]{C0C0C0} 
     \textit{\textbf{HERRAMIENTA}}                          & \textit{\textbf{NODO}} & \multicolumn{1}{l|}{\cellcolor[HTML]{C0C0C0}\textit{\textbf{PROCESO}}} & \multicolumn{1}{l|}{\cellcolor[HTML]{C0C0C0}\textit{\textbf{CPU}}} & \multicolumn{1}{l|}{\cellcolor[HTML]{C0C0C0}\textit{\textbf{GPU}}} \\ \hline
     \cellcolor[HTML]{EFEFEF}\textit{\textbf{Energy Meter}} & \textbf{X}             & \textbf{-}                                                             & \textbf{X}                                                         & \textbf{-}                                                         \\ \hline
     \cellcolor[HTML]{EFEFEF}\textit{\textbf{Kepler}}       & \textbf{X}             & \textbf{X}                                                             & \textbf{X}                                                         & \textbf{X}                                                         \\ \hline
     \cellcolor[HTML]{EFEFEF}\textit{\textbf{Scaphandre}}   & \textbf{X}             & \textbf{X}                                                             & \textbf{X}                                                         & \textbf{-}                                                         \\ \hline
     \cellcolor[HTML]{EFEFEF}\textit{\textbf{Monitor}}      & \textbf{X}             & \textbf{X}                                                             & \textbf{X}                                                         & \textbf{-}                                                         \\ \hline
     \end{tabular}
     \label{tab:arte-herramientas-capacidad-metricas}
     \caption{Capacidad de toma de metricas de las Herramientas}
\end{table}

\subsubsection*{Kepler}
 Kubernetes Efficient Power Level Exporter (Kepler) es una herramienta de monitorización de consumo de energía pensada para ser utilziada en entornos con multiples nodos gestionados por Kubernetes sobre los cuales se pretende medir el consumo energetico de una carga de trabajo desplegada sobre un contenedor, y saber cual es el consumo energetico que este presenta.\\

% DISTINGUE ENTRE:
%    - Consumo idel del sistema
%    - Consumo del sistema en proceso concreto
 Se basa en Intel RAPL para la obtencio de metricas de consumo del sistema y se apoya en nvidia-smi y sus librerias internar para obtener metricas de consumo de la GPU. Toma metricas cada 1 segundo, por defecto, y posteriormente aplica un post-procesamiento que resulta en un archivo compatible con Prometheus. Este es un software de monitorizacion de alertas compatible con Grafana, el cual no es necesario utilizar para tomar metricas con Kepler.\
 
 El conjunto de datos que proporciona Kepler en la salida se divide en tres bloques: El primero hace relacion al consumo energetico y metricas del programa, el segundo al consumo energetico a nivel de nodo y el tercero al consumo del proceso individual. Tanto en el segundo como el tercer bloque, distinguen metricas por zona de metricas(equivalentes a las que establece RAPL) y por core de la CPU.\

 Cabe destacar que pese a que Kepler es capaz de identificar que nodos han estado implicados en el proceso, no es capaz de diferenciar hacerca del porcentaje de implicacion que cada nodo a tenido en el proceso, reparte elconsumo energetico de forma equitativa.\\

 Los resultado que se obtiene del posprocesamiento con kepler, se dividen en dos tipos: Enrgia (Energy), representada en microjulios (μJ) de forma aculativa el consumo energetico del hardward; Y potencia (Power), representada en microwatios (μW) el consumo del hardware a lo largo del tiempo de medicion. 
 
 % Understanding Energy vs Power
 %    - Energy: Measured in microjoules (μJ) as cumulative counters from hardware
 %    - Power: Calculated as rate in microwatts (μW) using Power = ΔEnergy / Δtime

 %LIMITACION
 Segun se advierte en la documentacion oficial, Kepler presenta dificultades para medir de forma precisa el consumo energetico en entornos de carga mixta con patrones variables de cómputo y E/S, el modelo de atribución proporcional al tiempo de CPU puede ser inexacto, ya que no distingue entre el alto consumo de un proceso de cómputo intensivo y el menor gasto de uno que pasa gran parte del tiempo esperando operaciones de E/S, a pesar de que ambos puedan tener el mismo uso de CPU. Tambien presenta dificultades para realizar estimaciones sobre cargas de trabajo de alto rendimiento que usan instrucciones especializadas, como las vectoriales, la atribución basada únicamente en el tiempo de CPU no refleja el consumo energético real, que puede ser significativamente mayor. Por lo tanto, Kepler no es adecuado para decisiones de presupuesto energético absoluto ni para una optimización de grano fino que requiera mediciones precisas por proceso, ya que está diseñado para comparaciones relativas y análisis de tendencias, no para métricas exactas a nivel de proceso individual.

\subsubsection*{Scaphandre}


\subsubsection*{Energy Meter}


\subsubsection*{Monitor}